%% kapa: exam=2 round=27 year=2016 period=3 subject=compensation_law subject_ko=감정평가 및 보상법규 seq=4 number=4 points=10
%% source: Q-Net 감정평가사 기출문제 (cst00309) artlSeq=5185401 file=extracted/1008/5185401/제27회 감정평가사 2차시험 기출문제/제27회 감정평가사 2차 시험문제지/(3교시 시험문제지)감정평가 및 보상법규.pdf
%% conversion: kordoc 4.24.51 -> tools/gen2.py (auto); review=proofread-gpt-sol; flags=none
\begin{longquestion}{【문제 4】}
국토교통부장관은 감정평가업자 甲이「부동산 가격공시 및 감정평가에 관한 법률」(이하 ‘부동산공시법’이라 함) 제29조제2항에 따른 업무 범위를 위반하여 업무를 행하였다는 이유로 甲에게 3개월 업무정지처분을 하였다. 甲은 이러한 처분에 불복하여 취소소송을 제기하였으나 소송계속 중 3개월의 정지기간이 경과되었다. 부동산공시법 제38조제6항에 근거하여 제정된 부동산공시법 시행령 제77조 [별표 2] ‘감정평가업자의 설립인가의 취소와 업무의 정지에 관한 기준’에 따르면, 위 위반행위의 경우 위반횟수에 따라 가중처분을 하도록 규정하고 있다(1차 위반 시 업무정지 3개월, 2차 위반 시 업무정지 6개월, 3차 위반 시 업무정지 1년). 甲은 업무정지처분의 취소를 구할 법률상 이익이 있는가? (10점) \par
\end{longquestion}
